\chapter{VLANs and Access Ports}\label{ch:vlans}
\bp{2.2}

\begin{outcomes}
By the end of this chapter you will be able to:
\begin{tight}
  \item \textbf{Explain} what a VLAN changes about a switch, in terms of the
        learning-and-flooding algorithm from \chref{ethernet}.
  \item \textbf{Create} VLANs and assign access ports to them.
  \item \textbf{Verify} VLAN membership from three different angles, and know
        which one to trust when they disagree.
  \item \textbf{Diagnose} the classic VLAN faults: a port in the wrong VLAN, a
        VLAN that does not exist, and a VLAN that is shut down.
\end{tight}
\end{outcomes}

\begin{prereq}
\chref{ethernet} for the MAC table and flooding; \chref{cli} for moving around
the command line.
\end{prereq}

\begin{keyterms}
VLAN $\bullet$ broadcast domain $\bullet$ access port $\bullet$ VLAN database
$\bullet$ default VLAN $\bullet$ native VLAN $\bullet$ management VLAN $\bullet$
SVI $\bullet$ VLAN 1 $\bullet$ shutdown VLAN
\end{keyterms}

\section{What a VLAN actually does}

\begin{definitionbox}[VLAN]
A \term{VLAN} is a broadcast domain implemented in software on a switch. Ports in
VLAN 10 behave as though they were on a physically separate switch from ports in
VLAN 20: a broadcast from one never reaches the other, and a frame cannot cross
between them without passing through something that routes.
\end{definitionbox}

\begin{conceptbox}[One line changes in the switching algorithm]
Recall the algorithm from \chref{ethernet}: learn the source, then forward,
flood, or filter. VLANs change exactly one thing --- \textbf{every step is now
scoped to a VLAN}.

\begin{tight}
  \item Learning is per VLAN. The same MAC address may legitimately appear in two
        VLANs on the same switch.
  \item Flooding is per VLAN. An unknown unicast in VLAN 10 goes out every
        \emph{VLAN 10} port and no others.
  \item A frame never leaves the VLAN it arrived in. Not once, not for broadcast,
        not for anything.
\end{tight}
That is the whole mechanism. Everything else about VLANs --- trunking,
inter-VLAN routing, voice VLANs --- is machinery for coping with the consequences
of that one rule.
\end{conceptbox}

\begin{alertbox}[The consequence people forget]
If two devices are in different VLANs they cannot communicate, \emph{even if they
are in the same IP subnet, plugged into adjacent ports on the same switch}. The
switch will not carry the frame. Conversely two devices in the same VLAN but
different subnets cannot communicate either, because IP will not deliver it.

VLAN and subnet are separate ideas that must be kept aligned by design. The
convention that saves everyone trouble is one VLAN, one subnet, and a number that
matches: VLAN 10 is \cmd{192.168.10.0/24}, VLAN 20 is \cmd{192.168.20.0/24}.
\end{alertbox}

\section{Creating VLANs and assigning ports}

\begin{config}{three VLANs and the access ports that use them}
! Create the VLANs and name them. A name is documentation that
! survives on the device -- use it.
vlan 10
 name Staff
vlan 20
 name Guest
vlan 30
 name Voice
exit
!
! Assign a single access port.
interface GigabitEthernet0/5
 description Staff desk 5
 switchport mode access
 switchport access vlan 10
 no shutdown
!
! Assign a block of ports in one go.
interface range GigabitEthernet0/6 - 12
 description Staff desks 6-12
 switchport mode access
 switchport access vlan 10
!
interface range GigabitEthernet0/13 - 18
 description Guest desks
 switchport mode access
 switchport access vlan 20
\end{config}

\begin{conceptbox}[Why \cmd{switchport mode access} matters]
Without it the port is left in dynamic auto mode, which means it will negotiate
and may become a \emph{trunk} if the device on the other end asks for one. On an
access port that is both a fault waiting to happen and a security hole --- an
attacker who can negotiate a trunk can reach every VLAN.

Set it explicitly. It is one line, and it is the difference between a port that
does what you configured and a port that does what the far end suggests.
\end{conceptbox}

\section{Verifying, three ways}

\begin{verify}{SW1}{show vlan brief}
VLAN Name                             Status    Ports
---- -------------------------------- --------- -------------------------------
1    default                          active    Gi0/1, Gi0/2, Gi0/3, Gi0/4
10   Staff                            active    Gi0/5, Gi0/6, Gi0/7, Gi0/8
                                                Gi0/9, Gi0/10, Gi0/11, Gi0/12
20   Guest                            active    Gi0/13, Gi0/14, Gi0/15
30   Voice                            active
1002 fddi-default                     act/unsup
\end{verify}

This is the view by VLAN. Note that VLAN 30 exists but has no ports --- it was
created and never used, which is a common and harmless state, and also a common
symptom of a half-finished change.

\begin{verify}{SW1}{show interfaces GigabitEthernet0/5 switchport}
Name: Gi0/5
Switchport: Enabled
Administrative Mode: static access
Operational Mode: static access
Administrative Trunking Encapsulation: dot1q
Access Mode VLAN: 10 (Staff)
Trunking Native Mode VLAN: 1 (default)
Voice VLAN: none
\end{verify}

This is the view by port, and it is the one to trust when something is odd.
\textbf{Administrative} mode is what you configured; \textbf{operational} mode is
what the port has actually become. When those two differ, you have found your
fault.

\begin{verify}{SW1}{show mac address-table vlan 10}
          Mac Address Table
-------------------------------------------
Vlan    Mac Address       Type        Ports
----    -----------       --------    -----
  10    000a.f36e.7712    DYNAMIC     Gi0/5
  10    0060.5c2b.19a4    DYNAMIC     Gi0/7
\end{verify}

This is the view by evidence: which devices are actually present in the VLAN. If
a port is in VLAN 10 and no MAC address ever appears against it, nothing is
plugged in or the device is silent --- and that distinction is the start of most
access-layer troubleshooting.

\begin{pitfall}
\begin{tight}
  \item \textbf{Assigning a port to a VLAN that does not exist.} On many
        platforms the switch creates it silently; on others the port simply stops
        forwarding. Either way \cmd{show vlan brief} tells you within seconds.
  \item \textbf{Leaving everything in VLAN 1.} VLAN 1 cannot be deleted and
        carries control traffic. Do not put users on it, and do not use it as the
        management VLAN.
  \item \textbf{Forgetting that a VLAN can be shut down.} \cmd{shutdown} inside
        \cmd{vlan 10} stops the whole VLAN forwarding on that switch while every
        port still shows as correctly assigned. It is an unusual fault and
        therefore a nasty one.
  \item \textbf{Deleting a VLAN and leaving ports assigned to it.} Those ports
        become inactive. \cmd{show vlan brief} will not list them at all, which is
        confusing until you know to look for their absence.
\end{tight}
\end{pitfall}

\begin{breakfix}[A new desk cannot reach anything, and the port looks fine.]
A user on \cmd{Gi0/9} has no connectivity. The port is up, the cable is good, and
the configuration reads exactly like the working port next to it.

\begin{verify}{SW1}{show vlan brief | include Staff|Gi0/9}
10   Staff                            active    Gi0/5, Gi0/6, Gi0/7, Gi0/8
\end{verify}

\cmd{Gi0/9} is not listed against VLAN 10 --- nor against any other VLAN.

\begin{verify}{SW1}{show interfaces GigabitEthernet0/9 switchport}
Name: Gi0/9
Switchport: Enabled
Administrative Mode: static access
Operational Mode: static access
Access Mode VLAN: 15 (Inactive)
\end{verify}

\textbf{Diagnosis.} The port was assigned to VLAN 15, which does not exist on this
switch --- most likely a typing slip for 10, or a VLAN that was deleted after the
port was configured. The word \cmd{(Inactive)} is the whole diagnosis: the port is
administratively correct and operationally orphaned.

\textbf{Fix.} Either create VLAN 15 if it was meant to exist, or reassign the port:
\cmd{switchport access vlan 10}. Confirm with \cmd{show vlan brief} that
\cmd{Gi0/9} now appears in the Staff list, then confirm a MAC address learns
against it.
\end{breakfix}

\begin{lab}{Divide one switch into three networks}{Packet Tracer}
\textbf{Topology.} One switch, six PCs --- two per VLAN --- all in the same IP
subnet to begin with.

\begin{tasks}
  \item With every port in VLAN 1, confirm all six PCs can ping each other.
        Record the MAC address table.
  \item Create VLANs 10, 20 and 30 with names, and move two PCs into each.
  \item Re-test. Which pings now succeed and which fail? Explain each result in
        terms of broadcast domains, not IP.
  \item Run all three verification commands from this chapter and reconcile them.
  \item Assign \cmd{Gi0/11} to VLAN 99, which you have not created. Record what
        \cmd{show vlan brief} and \cmd{show interfaces switchport} each say.
  \item Enter \cmd{vlan 20} and type \cmd{shutdown}. Confirm the ports still look
        correctly assigned while nothing passes. Undo it.
  \item \textbf{Extension.} Re-address the PCs so each VLAN is its own subnet.
        Note that inter-VLAN pings still fail --- and that fixing it needs
        \chref{trunking}.
\end{tasks}
\end{lab}

\begin{checkpoint}
\begin{tightnum}
  \item Two PCs are in the same IP subnet, plugged into the same switch, in
        VLANs 10 and 20. Can they ping each other? Justify your answer at layer 2.
  \item Which command distinguishes what you configured on a port from what the
        port has actually become?
  \item A port shows \cmd{Access Mode VLAN: 15 (Inactive)}. State the two
        possible causes.
\end{tightnum}
\end{checkpoint}

\begin{chaptersummary}
\begin{tight}
  \item A VLAN is a broadcast domain in software. Learning, flooding and
        forwarding all become per-VLAN operations.
  \item A frame never leaves its VLAN without being routed.
  \item Keep VLAN and subnet aligned one-to-one; mismatches produce faults that
        look like everything and nothing.
  \item Always set \cmd{switchport mode access} explicitly, or the port may
        negotiate a trunk.
  \item Verify three ways: \cmd{show vlan brief} by VLAN,
        \cmd{show interfaces switchport} by port, and the MAC table by evidence.
        Administrative versus operational mode is where faults show.
\end{tight}
\end{chaptersummary}

\begin{reviewq}
  \item \textbf{(MCQ)} Which command assigns an access port to VLAN 10?
        (a)~\cmd{vlan 10} (b)~\cmd{switchport access vlan 10}
        (c)~\cmd{switchport trunk vlan 10} (d)~\cmd{ip vlan 10}
  \item \textbf{(MCQ)} A port reports \cmd{Administrative Mode: dynamic auto} and
        \cmd{Operational Mode: trunk}. What has happened?
        (a)~it was configured as a trunk (b)~it negotiated a trunk with the far
        end (c)~the VLAN is inactive (d)~the port is errdisabled
  \item \textbf{(Short)} Explain why leaving a port in dynamic auto is a security
        problem as well as an operational one.
  \item \textbf{(Short)} A VLAN exists, ports are assigned, and nothing forwards.
        Name a cause that \cmd{show vlan brief} would reveal and one it would not.
  \item \textbf{(Applied)} You are asked to segment a 24-port switch for staff,
        guests and printers, with guests unable to reach the other two. Write the
        configuration, and state what additional device or feature you would need
        for the isolation requirement to hold once routing is added.
\end{reviewq}
